%%%%%%%%%%%%%%%%%%%%%%%%%%%%%%%%%%%%%%%%%%%%%%%%%%%%%%%%%%%%%%%%%%%%%%%%%%%%%%%%
%%  content/02-background.tex — فصل ۲: مفاهیم پایه و مبانی نظری
%%%%%%%%%%%%%%%%%%%%%%%%%%%%%%%%%%%%%%%%%%%%%%%%%%%%%%%%%%%%%%%%%%%%%%%%%%%%%%%%
\section{مفاهیم پایه و مبانی نظری}

%%%%%%%%%%%%%%%%%%%%%%%%%%%%%%%%%%%%%%%%%%%%%%%%%%%%%%%%%%%%%%%%%%%%%%%%%%%%%%%%
\begin{frame}{بازیابی‌افزودهٔ تولید}{\en{Retrieval-Augmented Generation}}
  \begin{PTdef}[\en{RAG}]
    الگویی که پیش از تولید پاسخ، اسناد مرتبط را از یک منبع بیرونی بازیابی می‌کند
    و آن‌ها را همراه پرسش به مدل زبانی می‌سپارد؛ بنابراین دانش سامانه از
    \hl{وزن‌های مدل} جدا و در \hl{زمان اجرا} قابل به‌روزرسانی است.
  \end{PTdef}

  \begin{itemize}
    \item دانش تازه بدون آموزش دوباره افزوده می‌شود.
    \item هر جملهٔ پاسخ به سند مشخصی گره می‌خورد، پس \term{استنادپذیر} است.
    \item کیفیت پاسخ مستقیماً تابع کیفیت \hl{بازیابی} است، نه فقط اندازهٔ مدل.
  \end{itemize}
\end{frame}

%%%%%%%%%%%%%%%%%%%%%%%%%%%%%%%%%%%%%%%%%%%%%%%%%%%%%%%%%%%%%%%%%%%%%%%%%%%%%%%%
\begin{frame}{بازیابی معنایی}{از تطبیق واژه تا تطبیق معنا}
  پرسش و سند هر دو به برداری در فضایی چندهزار‌بعدی نگاشته می‌شوند و مرتبط‌بودن با
  \term{شباهت کسینوسی} سنجیده می‌شود:
  \[
    \operatorname{sim}(q,d)=
      \frac{\mathbf{v}_q\cdot\mathbf{v}_d}
           {\lVert\mathbf{v}_q\rVert\;\lVert\mathbf{v}_d\rVert}
  \]

  \begin{itemize}
    \item هم‌معنایی و صورت‌های واژگانی گوناگون به‌طور طبیعی پوشش داده می‌شوند.
    \item نگاشت یک‌بار محاسبه و نمایه می‌شود؛ در زمان پاسخ‌دهی فقط جست‌وجوی
          نزدیک‌ترین همسایه لازم است.
  \end{itemize}

  \begin{PTnote}
    بازیابی معنایی \hl{یادآوری} بالایی دارد اما ترتیبش خام است؛ همین شکاف، انگیزهٔ
    مرحلهٔ بازرتبه‌بندی در معماری پیشنهادی است.
  \end{PTnote}
\end{frame}

%%%%%%%%%%%%%%%%%%%%%%%%%%%%%%%%%%%%%%%%%%%%%%%%%%%%%%%%%%%%%%%%%%%%%%%%%%%%%%%%
\begin{frame}{بازرتبه‌بندی}{\en{Cross-Encoder Reranking}}
  \begin{itemize}
    \item بازیاب برداری پرسش و سند را \hl{جداگانه} می‌بیند؛ بازرتبه‌بند آن‌ها را
          \hl{با هم} به مدل می‌دهد و امتیاز مرتبط‌بودن مستقیم می‌گیرد.
    \item هزینهٔ محاسباتی بالاتر است، پس تنها روی \term{چند دهک بالای} نتایج
          بازیابی اجرا می‌شود.
  \end{itemize}

  \begin{PTresult}
    ترکیب «بازیابی گسترده + بازرتبه‌بندی باریک» هم دقت را بالا می‌برد و هم بافتار
    ورودی مدل زبانی را کوتاه می‌کند.
  \end{PTresult}
\end{frame}

%%%%%%%%%%%%%%%%%%%%%%%%%%%%%%%%%%%%%%%%%%%%%%%%%%%%%%%%%%%%%%%%%%%%%%%%%%%%%%%%
\begin{frame}{سنجه‌های ارزیابی}
  \begin{description}
    \item[\en{Precision@k}] نسبت اسناد مرتبط در \(k\) نتیجهٔ نخست.
    \item[\en{NDCG@k}] همان دقت، اما با وزن‌دهی به \hl{جایگاه} هر سند:
  \end{description}
  \[
    \operatorname{DCG}@k=\sum_{i=1}^{k}\frac{2^{r_i}-1}{\log_2(i+1)},
    \qquad
    \operatorname{NDCG}@k=\frac{\operatorname{DCG}@k}{\operatorname{IDCG}@k}
  \]
  \begin{itemize}
    \item \(r_i\) درجهٔ مرتبط‌بودن سند جایگاه \(i\) و \(\operatorname{IDCG}@k\)
          مقدار همان سنجه در ترتیب آرمانی است.
    \item مقدار ۱ یعنی ترتیب بازیابی‌شده با ترتیب آرمانی یکی است.
  \end{itemize}
\end{frame}
